\chapter{アルゴリズムの基礎}
\label{ch:algorithmic-foundations}

本章では，離散最適輸送問題を解くためのアルゴリズムを扱う．
第2章で定式化した Kantorovich 問題の離散版
\[
  \MKD_{\mathbf{C}}(\mathbf{a}, \mathbf{b})
  = \min_{\mathbf{P} \in \CouplingsD(\mathbf{a}, \mathbf{b})}
  \inner{\mathbf{C}}{\mathbf{P}}
\]
およびその双対問題を，組合せ最適化と線形計画法の手法で解くことが目標である．

これらのアルゴリズムの起源は第二次世界大戦前後にまで遡る．
Tolstoi (1930)，Hitchcock (1941)，Kantorovich (1942) による輸送問題の定式化に対し，
Dantzig の単体法（1947）が厳密な解法を初めて与えた．
最適輸送問題は線形計画の特殊なケースであるが，
最小コストフロー問題と等価であるという構造を持ち，
この構造を利用した効率的なアルゴリズムが開発されてきた．

%% ============================================================
%%  3.1 標準形への変換
%% ============================================================
\section{Kantorovich 問題の標準形}
\label{sec:standard-form}

Kantorovich 問題を線形計画の標準形に変換する．
行列 $\mathbf{P} \in \R^{n \times m}$ を列ごとに並べたベクトル
$\mathbf{p} \in \R^{nm}$（$\mathbf{p}$ の第 $i + n(j-1)$ 成分が $P_{i,j}$）を用いると，
周辺分布の制約は
\[
  \mathbf{A} \mathbf{p}
  = \begin{bmatrix} \mathbf{a} \\ \mathbf{b} \end{bmatrix},
  \quad
  \mathbf{A} =
  \begin{bmatrix}
    \ones_n^\top \otimes \Identity_m \\
    \Identity_n \otimes \ones_m^\top
  \end{bmatrix}
  \in \R^{(n+m) \times nm}
\]
と書ける．ここで $\otimes$ はクロネッカー積である．
コスト行列 $\mathbf{C}$ も同様にベクトル $\mathbf{c} \in \R^{nm}$ に並べると，
Kantorovich 問題は標準形の線形計画
\[
  \MKD_{\mathbf{C}}(\mathbf{a}, \mathbf{b})
  = \min_{\substack{\mathbf{p} \in \R_{\geq 0}^{nm} \\
  \mathbf{A}\mathbf{p} = [\mathbf{a}^\top, \mathbf{b}^\top]^\top}}
  \mathbf{c}^\top \mathbf{p}
\]
と表される．

\begin{definition}{クロネッカー積}{kronecker}
  行列 $\mathbf{A} \in \R^{p \times q}$ と $\mathbf{B} \in \R^{r \times s}$ の
  \textbf{クロネッカー積}（Kronecker product）$\mathbf{A} \otimes \mathbf{B} \in \R^{pr \times qs}$
  を，$(i, j)$ ブロックが $A_{i,j} \mathbf{B}$ であるブロック行列として定義する：
  \[
    \mathbf{A} \otimes \mathbf{B} =
    \begin{pmatrix}
      A_{1,1}\mathbf{B} & \cdots & A_{1,q}\mathbf{B} \\
      \vdots & \ddots & \vdots \\
      A_{p,1}\mathbf{B} & \cdots & A_{p,q}\mathbf{B}
    \end{pmatrix}.
  \]
\end{definition}

\begin{remark}{制約の冗長性}{constraint-redundancy}
  行列 $\mathbf{A}$ の $n + m$ 本の行のうち1本は冗長である．
  実際，最初の $n$ 行の和と残りの $m$ 行の和はいずれも $\ones_{nm}^\top$ に等しく，
  これは $\sum_i a_i = \sum_j b_j = 1$ という条件に対応する．
  したがって $\mathbf{A}$ の階数は $n + m - 1$ であり，
  最適解を探索すべき $\CouplingsD(\mathbf{a}, \mathbf{b})$ の頂点は
  最大 $n + m - 1$ 個の非零成分を持つ（命題~\ref{prop:extremal-tree}）．
\end{remark}

%% ============================================================
%%  3.2 離散 C-変換
%% ============================================================
\section{離散 $C$-変換}
\label{sec:discrete-c-transform}

第2章で連続版の $c$-変換（定義~\ref{def:c-transform}）を導入した．
ここではその離散版を定義し，双対問題を1変数の最適化に帰着させる性質を示す．

\begin{definition}{離散 $C$-変換}{discrete-c-transform}
  コスト行列 $\mathbf{C} \in \R^{n \times m}$ とベクトル $\mathbf{f} \in \R^n$ に対して，
  $\mathbf{f}$ の\textbf{$C$-変換} $\mathbf{f}^{\,C} \in \R^m$ を
  \[
    (\mathbf{f}^{\,C})_j \defeq \min_{i \in \range{n}} (C_{i,j} - f_i)
  \]
  で定義する．同様に，$\mathbf{g} \in \R^m$ の $\bar{C}$-変換 $\mathbf{g}^{\,\bar{C}} \in \R^n$ を
  \[
    (\mathbf{g}^{\,\bar{C}})_i \defeq \min_{j \in \range{m}} (C_{i,j} - g_j)
  \]
  で定義する．
\end{definition}

$\mathbf{f}$ を固定したとき，$f_i + g_j \leq C_{i,j}$（$\forall\, i, j$）を満たす
$\mathbf{g}$ のうち成分ごとに最大のものが $\mathbf{f}^{\,C}$ である．
実際，$f_i + g_j \leq C_{i,j}$ がすべての $i$ で成り立つことは
$g_j \leq \min_i (C_{i,j} - f_i) = (\mathbf{f}^{\,C})_j$ と同値であるから，
制約を満たす最大の $g_j$ は $(\mathbf{f}^{\,C})_j$ に他ならない．
さらに，任意の実行可能な $(\mathbf{f}, \mathbf{g})$ に対して
$\mathbf{g} \leq \mathbf{f}^{\,C}$ であるから
$\inner{\mathbf{g}}{\mathbf{b}} \leq \inner{\mathbf{f}^{\,C}}{\mathbf{b}}$
（$\mathbf{b} \geq 0$ による）が成り立つ．
したがって，双対問題は $\mathbf{f}$ のみの最大化問題に帰着される：
\[
  \MKD_{\mathbf{C}}(\mathbf{a}, \mathbf{b})
  = \max_{\mathbf{f} \in \R^n}
  \left[ \inner{\mathbf{f}}{\mathbf{a}} + \inner{\mathbf{f}^{\,C}}{\mathbf{b}} \right].
\]
この定式化は\textbf{半双対}（semi-dual）と呼ばれる．

$C$-変換と $\bar{C}$-変換の交互適用には以下の性質がある．

\begin{proposition}{$C$-変換の性質}{c-transform-properties}
  以下が成り立つ（不等号はすべて成分ごと）：
  \begin{enumerate}
    \item \textbf{単調性の反転}：$\mathbf{f} \leq \mathbf{f}' \implies \mathbf{f}^{\,C} \geq \mathbf{f}'^{\,C}$．
    \item \textbf{拡大性}：$\mathbf{f}^{\,C\bar{C}} \geq \mathbf{f}$，
      $\mathbf{g}^{\,\bar{C}C} \geq \mathbf{g}$．
    \item \textbf{冪等性}：$\mathbf{f}^{\,C\bar{C}C} = \mathbf{f}^{\,C}$．
  \end{enumerate}
\end{proposition}

\begin{proof}
  \textbf{(1)}
  $\mathbf{f} \leq \mathbf{f}'$ ならば各 $j$ に対して
  $C_{i,j} - f_i \geq C_{i,j} - f'_i$ であるから，
  $\min_i (C_{i,j} - f_i) \geq \min_i (C_{i,j} - f'_i)$，
  すなわち $(\mathbf{f}^{\,C})_j \geq (\mathbf{f}'^{\,C})_j$ が成り立つ．

  \textbf{(2)}
  $C\bar{C}$-変換の定義を展開すると，
  \[
    (\mathbf{f}^{\,C\bar{C}})_i
    = \min_{j} \left( C_{i,j} - (\mathbf{f}^{\,C})_j \right)
    = \min_{j} \left( C_{i,j} - \min_{i'} (C_{i',j} - f_{i'}) \right).
  \]
  $-\min_{i'} (C_{i',j} - f_{i'}) \geq -(C_{i,j} - f_i) = f_i - C_{i,j}$ であるから，
  \[
    (\mathbf{f}^{\,C\bar{C}})_i \geq \min_{j} (C_{i,j} + f_i - C_{i,j}) = f_i.
  \]
  $\mathbf{g}^{\,\bar{C}C} \geq \mathbf{g}$ も同様．

  \textbf{(3)}
  $\mathbf{g} = \mathbf{f}^{\,C}$ とおくと，(2) より
  $\mathbf{g}^{\,\bar{C}} = \mathbf{f}^{\,C\bar{C}} \geq \mathbf{f}$ が成り立つ．
  (1) を適用すると $\mathbf{f}^{\,C\bar{C}C} \leq \mathbf{f}^{\,C}$ が得られる．
  一方，(2) より $\mathbf{f}^{\,C\bar{C}C} \geq \mathbf{f}^{\,C}$ であるから，
  $\mathbf{f}^{\,C\bar{C}C} = \mathbf{f}^{\,C}$ が成り立つ．
\end{proof}

\begin{remark}{}{c-transform-plateau}
  性質 (3) は，$C$-変換と $\bar{C}$-変換を交互に適用しても
  2回目以降は改善が得られないことを意味する．
  すなわち，$\mathbf{f} \to \mathbf{f}^{\,C} \to \mathbf{f}^{\,C\bar{C}}
  \to \mathbf{f}^{\,C\bar{C}C} = \mathbf{f}^{\,C} \to \cdots$
  と停滞する．
  したがって，$C$-変換の交互適用だけでは最適解には到達できず，
  ネットワーク単体法やオークションアルゴリズムなどの手法が必要となる．
\end{remark}

%% ============================================================
%%  3.3 輸送多面体の頂点
%% ============================================================
\section{輸送多面体の頂点}
\label{sec:vertices}

線形計画の最適解は実行可能領域の\textbf{頂点}（端点）で達成される．
$\CouplingsD(\mathbf{a}, \mathbf{b})$ の頂点の構造を理解することは，
アルゴリズムの設計に不可欠である．

\begin{definition}{二部グラフ}{bipartite-graph}
  ソース節点の集合 $V = \{1, \ldots, n\}$ と
  ターゲット節点の集合 $V' = \{1', \ldots, m'\}$ に対して，
  $V$ と $V'$ の間のすべての辺 $\{(i, j') \mid i \in \range{n},\, j \in \range{m}\}$ から成る
  グラフを\textbf{完全二部グラフ}（complete bipartite graph）と呼ぶ．
  輸送計画 $\mathbf{P}$ の\textbf{台グラフ}（support graph）を，
  $P_{i,j} > 0$ なる辺 $(i, j')$ から成るグラフ
  $G(\mathbf{P}) = (V \cup V',\, S(\mathbf{P}))$ で定義する．
\end{definition}

\begin{definition}{閉路}{cycle}
  グラフにおいて，相異なる節点 $v_1, v_2, \ldots, v_k$ と辺
  $(v_1, v_2), (v_2, v_3), \ldots, (v_{k-1}, v_k), (v_k, v_1)$ から成る経路を
  \textbf{閉路}（cycle）と呼ぶ．閉路を持たないグラフを\textbf{森}（forest）といい，
  連結な森を\textbf{木}（tree）という．
\end{definition}

\begin{proposition}{頂点の木構造}{extremal-tree}
  $\mathbf{P}$ が $\CouplingsD(\mathbf{a}, \mathbf{b})$ の頂点ならば，
  台グラフ $G(\mathbf{P})$ は閉路を持たない（すなわち森である）．
  特に，$\mathbf{P}$ の非零成分の数は $n + m - 1$ 以下である．
\end{proposition}

\begin{proof}
  背理法による．$\mathbf{P}$ が頂点であり，かつ $G(\mathbf{P})$ が
  閉路 $(i_1, j_1'), (i_2, j_1'), (i_2, j_2'), \ldots, (i_k, j_k'), (i_1, j_k')$ を
  含むと仮定する．

  閉路に沿って交互に $+\varepsilon$，$-\varepsilon$ の摂動を加えた行列
  $\mathbf{E}$ を構成する．
  $\varepsilon > 0$ を閉路上の $\mathbf{P}$ の最小値より小さく取れば，
  $\mathbf{Q} = \mathbf{P} + \mathbf{E}$ および $\mathbf{R} = \mathbf{P} - \mathbf{E}$ は
  いずれも非負である．
  $\mathbf{E}$ の各行和・各列和は $0$ であるから，
  $\mathbf{Q}, \mathbf{R} \in \CouplingsD(\mathbf{a}, \mathbf{b})$ である．
  $\mathbf{P} = (\mathbf{Q} + \mathbf{R}) / 2$ かつ $\mathbf{Q} \neq \mathbf{P}$ であるから，
  $\mathbf{P}$ が頂点であることに矛盾する．

  $k$ 個の節点を持つ閉路のないグラフの辺数は $k - 1$ 以下であるから，
  $n + m$ 個の節点を持つ $G(\mathbf{P})$ の辺数（$= \mathbf{P}$ の非零成分数）は
  $n + m - 1$ 以下である．
\end{proof}

%% ============================================================
%%  3.3 北西隅法
%% ============================================================
\section{北西隅法}
\label{sec:northwest}

ネットワーク単体法などの反復法を初期化するには，
$\CouplingsD(\mathbf{a}, \mathbf{b})$ の頂点を1つ求める必要がある．
\textbf{北西隅法}（north-west corner rule）は，高々 $n + m$ 回の演算で
頂点を構成する単純なヒューリスティックである．

\begin{algorithm}{alg:northwest}
  \textbf{Input:} $\mathbf{a} \in \simplex_n$, $\mathbf{b} \in \simplex_m$ \\
  \textbf{Output:} $\CouplingsD(\mathbf{a}, \mathbf{b})$ の頂点 $\mathbf{P}$ \\[4pt]
  1. $i \leftarrow 1$, $j \leftarrow 1$, $r \leftarrow a_1$, $c \leftarrow b_1$ \\
  2. \textbf{while} $i \leq n$ かつ $j \leq m$ \textbf{do}: \\
  3. \quad $t \leftarrow \min(r, c)$, \quad $P_{i,j} \leftarrow t$ \\
  4. \quad $r \leftarrow r - t$, \quad $c \leftarrow c - t$ \\
  5. \quad \textbf{if} $r = 0$: $i \leftarrow i + 1$, $r \leftarrow a_i$（$i \leq n$ のとき） \\
  6. \quad \textbf{if} $c = 0$: $j \leftarrow j + 1$, $c \leftarrow b_j$（$j \leq m$ のとき）
\end{algorithm}

\begin{example}{北西隅法の実行例}{northwest-example}
  $\mathbf{a} = (0.2, 0.5, 0.3)$，$\mathbf{b} = (0.5, 0.1, 0.4)$ のとき：
  \[
    \begin{pmatrix} 0.2 & 0 & 0 \\ 0.3 & 0.1 & 0.1 \\ 0 & 0 & 0.3 \end{pmatrix}.
  \]
  非零成分は $5 = 3 + 3 - 1$ 個であり，
  台グラフは閉路を持たないので，これは $\CouplingsD(\mathbf{a}, \mathbf{b})$ の頂点である．
\end{example}

%% ============================================================
%%  3.4 ネットワーク単体法
%% ============================================================
\section{ネットワーク単体法}
\label{sec:network-simplex}

\textbf{ネットワーク単体法}（network simplex）は，
最適輸送問題に特化した単体法の変種であり，
輸送多面体の頂点間を移動しながら最適解を探索する．
以下の3つの操作を反復する．

\subsection{相補的な双対変数の構成}
\label{subsec:dual-from-primal}

$\CouplingsD(\mathbf{a}, \mathbf{b})$ の頂点 $\mathbf{P}$ が与えられたとき，
$\mathbf{P}$ と\textbf{相補的}（complementary）な双対変数 $(\mathbf{f}, \mathbf{g})$ を求める．
すなわち，$P_{i,j} > 0$ なるすべての $(i, j)$ に対して $f_i + g_j = C_{i,j}$ を満たす
$(\mathbf{f}, \mathbf{g})$ を構成する．

命題~\ref{prop:extremal-tree} から，$G(\mathbf{P})$ は森であり辺数は $n + m - 1$ 以下である．
各連結成分（木）において，1つの節点の双対変数を $0$ に設定し，
木の幅優先探索（または深さ優先探索）により残りの変数を
$f_i + g_j = C_{i,j}$ から逐次的に定める．

\subsection{最適性の判定と更新}

得られた $(\mathbf{f}, \mathbf{g})$ が\textbf{双対実行可能}（$f_i + g_j \leq C_{i,j}$, $\forall i, j$）
であれば，主双対最適性条件（命題~\ref{prop:primal-dual-optimality}）から
$\mathbf{P}$ は最適であり，アルゴリズムは終了する．

双対実行可能でない場合，$f_i + g_j > C_{i,j}$ となる辺 $(i, j')$ が存在する．
この辺を $G(\mathbf{P})$ に追加すると，以下の2つの場合が起こりうる：

\begin{enumerate}
  \item \textbf{閉路が生じない場合}：
    $G$ は依然として森であり，$(\mathbf{f}, \mathbf{g})$ を再計算する．
    $\mathbf{P}$ は変更されない（退化ステップ）．
  \item \textbf{閉路が生じる場合}：
    閉路に沿って流量を調整する．
    閉路上の「正方向」の辺では $P_{i,j}$ を $\theta$ だけ増加させ，
    「負方向」の辺では $\theta$ だけ減少させる．
    $\theta$ は負方向の辺の最小流量に等しく，
    流量が $0$ になった辺を $G$ から除去する．
\end{enumerate}

\subsection{コストの改善}

閉路が生じた場合にコストが改善されることを示す．
閉路を $(i_1, j_1'), (j_1', i_2), (i_2, j_2'), \ldots, (i_l, j_l'), (j_l', i_1)$ と書き，
$i_1 = i$，$j_1 = j$（追加した辺）とする．
更新後の $\tilde{\mathbf{P}}$ は正方向の辺で $\theta$ 増加，負方向の辺で $\theta$ 減少するから，
\[
  \inner{\tilde{\mathbf{P}}}{\mathbf{C}} - \inner{\mathbf{P}}{\mathbf{C}}
  = \theta \left( \sum_{k=1}^l C_{i_k, j_k} - \sum_{k=1}^l C_{i_{k+1}, j_k} \right).
\]
$(\mathbf{f}, \mathbf{g})$ は $\mathbf{P}$ に相補的であるから，
既存の辺 $(i_k, j_k')$（$k \geq 2$）と $(i_{k+1}, j_k')$ について
$f_{i_k} + g_{j_k} = C_{i_k, j_k}$ および $f_{i_{k+1}} + g_{j_k} = C_{i_{k+1}, j_k}$ が成り立つ．
これらを代入すると，差は打ち消し合い
\[
  \sum_{k=1}^l C_{i_k, j_k} - \sum_{k=1}^l C_{i_{k+1}, j_k} = C_{i,j} - (f_i + g_j)
\]
が得られる．$(i, j)$ は $f_i + g_j > C_{i,j}$ なる辺として選ばれたから，
$C_{i,j} - (f_i + g_j) < 0$ であり，$\theta > 0$ ならばコストは厳密に減少する．

\begin{remark}{計算量}{network-simplex-complexity}
  ネットワーク単体法は多項式時間で停止することが Orlin (1997) により証明された．
  Tarjan (1997) は効率的なデータ構造を用いて
  $O\bigl((n+m)nm \log(n+m) \log((n+m)\norm{\mathbf{C}}_\infty)\bigr)$
  の計算量を達成した．
\end{remark}

%% ============================================================
%%  3.5 オークションアルゴリズム
%% ============================================================
\section{オークションアルゴリズム}
\label{sec:auction}

\textbf{オークションアルゴリズム}（auction algorithm）は
Bertsekas (1981) により提案された，
最適割当問題（$n = m$，均一重み）に特化した双対的手法である．

\subsection{$\varepsilon$-相補性条件}

最適割当問題（設定~\ref{set:assignment}）の最適解 $\sigma^*$ と
最適双対変数 $\mathbf{g}^*$ は
\[
  C_{i, \sigma^*(i)} - g^*_{\sigma^*(i)} = \min_{j} (C_{i,j} - g^*_j)
  \qquad (\forall\, i \in \range{n})
\]
を満たす（$c$-変換と相補性条件から従う）．

オークションアルゴリズムはこの条件の\textbf{近似版}を用いる．

\begin{definition}{$\varepsilon$-相補性条件}{epsilon-cs}
  部分割当 $(S, \xi)$（$S \subset \range{n}$，$\xi \colon S \to \range{n}$ は単射）と
  双対変数 $\mathbf{g} \in \R^n$ が
  \textbf{$\varepsilon$-相補性条件}（$\varepsilon$-complementary slackness; $\varepsilon$-CS）
  を満たすとは，すべての $i \in S$ に対して
  \[
    C_{i, \xi(i)} - g_{\xi(i)} \leq \varepsilon + \min_{j} (C_{i,j} - g_j)
  \]
  が成り立つことをいう．
\end{definition}

\subsection{アルゴリズムの更新規則}

未割当の $i \notin S$ に対して，調整コスト $C_{i,j} - g_j$ の
最小値と第2最小値を達成する添字を
\[
  j_i^1 \in \argmin_{j} (C_{i,j} - g_j), \qquad
  j_i^2 \in \argmin_{j \neq j_i^1} (C_{i,j} - g_j)
\]
とする．以下の更新を行う：

\begin{enumerate}
  \item \textbf{$\mathbf{g}$ の更新}：
    \[
      g_{j_i^1} \leftarrow g_{j_i^1}
      - \Bigl[\bigl(C_{i, j_i^2} - g_{j_i^2}\bigr) - \bigl(C_{i, j_i^1} - g_{j_i^1}\bigr) + \varepsilon\Bigr].
    \]
    括弧内は第2最小値と最小値の差に $\varepsilon$ を加えたものであり，常に $\geq \varepsilon > 0$ である．
  \item \textbf{$S$ と $\xi$ の更新}：
    $j_i^1$ に既に割当てられている $i'$（$\xi(i') = j_i^1$）があれば
    $S$ から $i'$ を除去する．$\xi(i) = j_i^1$ として $i$ を $S$ に追加する．
\end{enumerate}

\begin{proposition}{$\varepsilon$-CS の保存}{auction-eps-cs}
  各反復後も $\varepsilon$-CS 条件は保たれる．
\end{proposition}

\begin{proof}
  更新前に $\varepsilon$-CS が成り立つとする．
  更新後の双対変数を $\mathbf{g}^{\mathrm{n}}$ と書く．
  $\mathbf{g}^{\mathrm{n}}$ は $\mathbf{g}$ から $j_i^1$ 成分のみ減少させたものであるから
  $\mathbf{g}^{\mathrm{n}} \leq \mathbf{g}$（成分ごと）が成り立つ．

  新たに追加された $i$ について，更新の定義から
  $C_{i, j_i^1} - g^{\mathrm{n}}_{j_i^1} = \varepsilon + (C_{i, j_i^2} - g_{j_i^2})
  = \varepsilon + \min_{j \neq j_i^1}(C_{i,j} - g_j)$ であるから，
  $\mathbf{g}^{\mathrm{n}} \leq \mathbf{g}$ より
  $C_{i, j_i^1} - g^{\mathrm{n}}_{j_i^1}
  \leq \varepsilon + \min_j (C_{i,j} - g^{\mathrm{n}}_j)$ が成り立つ．

  $i' \neq i$（$i' \in S$）について，割当は変わらず $\mathbf{g}^{\mathrm{n}} \leq \mathbf{g}$ であるから
  $C_{i', \xi(i')} - g^{\mathrm{n}}_{\xi(i')}
  = C_{i', \xi(i')} - g_{\xi(i')}
  \leq \varepsilon + \min_j (C_{i',j} - g_j)
  \leq \varepsilon + \min_j (C_{i',j} - g^{\mathrm{n}}_j)$ が成り立つ．
\end{proof}

\begin{proposition}{反復回数の上界}{auction-iterations}
  オークションアルゴリズムは最大 $n \norm{\mathbf{C}}_\infty / \varepsilon$ 回の反復で停止する．
\end{proposition}

\begin{proof}
  アルゴリズムが $T > n\norm{\mathbf{C}}_\infty / \varepsilon$ 回反復して停止しないと仮定する．
  すると，$\xi$ の像に入っていない添字 $j$（一度も割当先にならなかった）が存在し，
  $g_j = 0$ のままである．
  一方，ある添字 $j'$ が $n$ 回以上更新されたとすると，
  各更新で $g_{j'}$ は少なくとも $\varepsilon$ 減少するから
  $g_{j'} \leq -n\varepsilon < -\norm{\mathbf{C}}_\infty$ となる．
  このとき，任意の $i$ に対して
  $C_{i,j'} - g_{j'} > \norm{\mathbf{C}}_\infty + \norm{\mathbf{C}}_\infty
  \geq C_{i,j} + \norm{\mathbf{C}}_\infty > C_{i,j} - g_j$ であるから，
  $j'$ が $j_i^1 = \argmin_j (C_{i,j} - g_j)$ となることはない．
  これは $j'$ がさらに更新されることに矛盾する．
  したがって，各添字の更新回数は $\norm{\mathbf{C}}_\infty / \varepsilon$ 以下であり，
  全体の反復回数は $n\norm{\mathbf{C}}_\infty / \varepsilon$ 以下である．
\end{proof}

\begin{proposition}{近似最適性}{auction-suboptimality}
  オークションアルゴリズムの出力する割当のコストは，
  最適値に対して $n\varepsilon$ 以内の近似解である．
\end{proposition}

\begin{proof}
  最適割当を $\sigma^*$，最適双対変数を $\mathbf{g}^*$ とし，最適値を
  $t^* = \sum_i C_{i,\sigma^*(i)} = \sum_i \min_j (C_{i,j} - g^*_j) + \sum_j g^*_j$
  と書く．
  アルゴリズム終了時の割当 $\xi$ と双対変数 $\mathbf{g}$ に対して，
  $\mathbf{g}$ の双対実行可能性（$\mathbf{g} \leq \mathbf{g}^*$ の意味での劣最適性）から
  \[
    t^* \geq \sum_i \min_j (C_{i,j} - g_j) + \sum_j g_j.
  \]
  $\varepsilon$-CS 条件 $C_{i,\xi(i)} - g_{\xi(i)} \leq \varepsilon + \min_j (C_{i,j} - g_j)$ を用いると，
  \[
    \sum_i \min_j (C_{i,j} - g_j)
    \geq \sum_i \left( C_{i,\xi(i)} - g_{\xi(i)} - \varepsilon \right).
  \]
  $\xi$ は全単射であるから $\sum_i g_{\xi(i)} = \sum_j g_j$ が成り立ち，
  \[
    t^* \geq \sum_i C_{i,\xi(i)} - n\varepsilon.
  \]
  一方 $\sum_i C_{i,\xi(i)} \geq t^*$（$\xi$ は一般に最適でないから）より，
  $t^* \leq \sum_i C_{i,\xi(i)} \leq t^* + n\varepsilon$ が得られる．
\end{proof}

\begin{remark}{$\varepsilon$-スケーリング}{epsilon-scaling}
  $\varepsilon$ を徐々に減少させる\textbf{$\varepsilon$-スケーリング}を用いると，
  計算量は $O(n^3 \log \norm{\mathbf{C}}_\infty)$ に改善される．
\end{remark}

%% ============================================================
%%  3.6 双対上昇法
%% ============================================================
\section{双対上昇法}
\label{sec:dual-ascent}

ネットワーク単体法が主問題の頂点を更新するのに対し，
\textbf{双対上昇法}（dual ascent method）は双対変数 $(\mathbf{f}, \mathbf{g})$ を
直接改善する手法であり，その起源は Borchardt (1865) にまで遡る．
K\"{o}nig--Egerv\'{a}ry の定理や Kuhn (1955) による\textbf{ハンガリアン法}
（Hungarian algorithm）はこの枠組みの代表例である．

\subsection{均衡辺と不活性辺}

双対変数 $(\mathbf{f}, \mathbf{g}) \in \PotentialsD(\mathbf{C})$ が与えられたとき，
辺 $(i, j)$ が\textbf{均衡}（balanced）であるとは $f_i + g_j = C_{i,j}$ が成り立つことをいい，
$f_i + g_j < C_{i,j}$ のとき\textbf{不活性}（inactive）であるという．
均衡辺から成る部分グラフを $G_{\mathbf{f}, \mathbf{g}}$ と記す．

\subsection{双対上昇ステップ}

$(\mathbf{f}, \mathbf{g})$ が最適でない場合，
双対目的関数 $\inner{\mathbf{f}}{\mathbf{a}} + \inner{\mathbf{g}}{\mathbf{b}}$ を
改善する方向が存在する．

集合 $S \subset \range{n}$，$S' \subset \range{m}$ を，
均衡辺グラフ $G_{\mathbf{f}, \mathbf{g}}$ における\textbf{到達可能性}から定める．
具体的には，$S$ 内のソース節点 $i$ から出る均衡辺の
ターゲット側の節点がすべて $S'$ に含まれるように $S, S'$ を選ぶ．
このとき，$\varepsilon > 0$ を十分小さく取れば，摂動した双対変数
\[
  (\tilde{\mathbf{f}}, \tilde{\mathbf{g}})
  = (\mathbf{f}, \mathbf{g})
  + \varepsilon(\ones_S, -\ones_{S'})
\]
は双対実行可能である．ここで $(\ones_S)_i = 1$（$i \in S$），$0$（$i \notin S$）である．
双対目的関数の変化量は
\[
  \varepsilon \left( \sum_{i \in S} a_i - \sum_{j \in S'} b_j \right)
\]
であり，$\sum_{i \in S} a_i > \sum_{j \in S'} b_j$ となる $S, S'$ を選べば改善が得られる．
このような $S, S'$ は，均衡辺グラフ上の最大フロー（Ford--Fulkerson アルゴリズム等）から
構成できる．

\begin{remark}{ネットワーク単体法との比較}{dual-ascent-vs-simplex}
  ネットワーク単体法は主問題の\textbf{実行可能}な頂点を維持し，
  最適性条件（双対実行可能性）は収束時にのみ達成される．
  双対上昇法は逆に，双対実行可能性を維持しつつ，
  主問題の実行可能性（周辺分布制約）は収束時にのみ達成される．
  割当問題（$n = m$，均一重み）の場合，双対上昇法は
  ハンガリアン法に帰着する．
\end{remark}
